\section{Restoration of Mesolimbic Reward Capacity as the Primary Treatment Target in Stimulant Use Disorder and Dopamine-Deficit Conditions}\label{18_restoration-restoration-of-mesolimbic-reward-capacity-as-the-primary-treatment-target-in-stimulant-use-disorder-and-dopamine-deficit-conditions}

\subsection{Abstract}\label{18_restoration-abstract}

Current treatment paradigms for stimulant use disorder prioritize abstinence, relapse prevention, and craving suppression as primary endpoints. This position paper argues that these are proxy measures that systematically misdirect both research and clinical practice. The correct treatment target is restoration of mesolimbic reward capacity --- operationally defined as the ability of natural reinforcers to generate sufficient dopaminergic signal to sustain motivated behavior and make continued existence feel worthwhile. We present evidence that (1) abstinence without reward restoration produces a structurally inevitable relapse pathway termed ``the anhedonic trap,'' (2) current proxy endpoints fail to measure the neurobiological substrate of the disorder, (3) existing interventions actively deepen dopaminergic deficits through stress-mediated mechanisms, and (4) a treatment framework oriented toward reward restoration requires different outcome measures, intervention strategies, and maintenance paradigms. The failure to measure hedonic capacity and motivational function in addiction trials represents not merely a methodological oversight but a conceptual error that has shaped decades of ineffective treatment. \textbf{Keywords:} stimulant use disorder; dopamine; D2 receptors; anhedonia; reward; mesolimbic; treatment outcomes; harm reduction

\subsection{}\label{18_restoration-section}

\subsection{}\label{18_restoration-section-1}

\subsection{}\label{18_restoration-section-2}

\subsection{}\label{18_restoration-section-3}

\subsection{}\label{18_restoration-section-4}

\subsection{Section 1: The Wrong Endpoints and What They Cost}\label{18_restoration-section-1-the-wrong-endpoints-and-what-they-cost}

\subsubsection{The Conceptual Separation of Three Domains}\label{18_restoration-the-conceptual-separation-of-three-domains}

Contemporary addiction treatment operates across three domains that are frequently conflated but represent fundamentally distinct phenomena: abstinence (behavioral), craving suppression (phenomenological), and reward restoration (neurobiological). Abstinence denotes the observable cessation of drug-seeking and drug-taking behavior, quantifiable through urine drug screens, timeline follow-back interviews, or treatment retention. Craving suppression refers to the subjective diminishment of drug-related urges, typically assessed through self-report instruments such as the Cocaine Craving Questionnaire or visual analog scales. Reward restoration --- the capacity of natural reinforcers to activate mesolimbic dopaminergic circuits at sufficient magnitude and reliability to sustain goal-directed behavior --- is a neurobiological state with measurable correlates in D2/D3 receptor availability, dopamine synthesis capacity, and functional connectivity of reward circuits. The field has treated these three domains as if they were causally linked: suppress craving, achieve abstinence, implicitly assume reward function returns. This assumption is mechanistically unfounded and empirically false.

\subsubsection{The Anhedonic Trap}\label{18_restoration-the-anhedonic-trap}

Abstinence without reward restoration produces what we term the anhedonic trap: a state in which the person is technically drug-free but experiences no alternative sources of positive valence, making relapse a structural inevitability rather than a moral failure. The neurobiological substrate of this trap is well-characterized. Chronic stimulant use produces D2/D3 receptor downregulation through GASP-1-mediated lysosomal degradation {[}{[}Bartlett et al., 2005, PNAS 102:11521--11526{]} (https://pmc.ncbi.nlm.nih.gov/articles/PMC1183554/){]}, ΔFosB/HDAC1-mediated transcriptional repression at the DRD2 promoter {[}\href{https://pubmed.ncbi.nlm.nih.gov/18632938/}{Renthal et al., 2008, J Neurosci 28:7344--7349}{]}, and epigenetic silencing via H3K27me3 accumulation {[}\href{https://pubmed.ncbi.nlm.nih.gov/28421257/}{Montalvo-Ortiz et al., 2017, Psychopharmacology 234:2385-- 2398}{]} and DNA methylation {[}{[}Nohesara et al., 2016, Am J Med Genet Part B 171:1180--1189{]} (https://pmc.ncbi.nlm.nih.gov/articles/PMC7115129/){]}. The only longitudinal PET study measuring both dopamine transporter and D2/D3 receptors found that while DAT showed approximately 20\% recovery at 9 months of abstinence, D2/D3 receptor availability showed no measurable change whatsoever {[}{[}Volkow et al., 2015, NeuroImage 121:20--28{]} (https://www.sciencedirect.com/science/article/abs/pii/S1053811915006461){]}. Robertson et al.~(2016) stated this explicitly: ``There are no published reports of recovery of D2/D3 receptors with drug abstinence in human subjects'' {[}{[}Robertson et al., 2016, Neuropsychopharmacology 41:1629--1636{]} (https://pmc.ncbi.nlm.nih.gov/articles/PMC4832026/){]}. When D2/D3 receptor density remains depressed, the dopaminergic signal generated by natural rewards --- food, social connection, work accomplishment, novelty, goal pursuit --- falls below the threshold required to encode incentive salience {[}{[}Berridge \& Robinson, 1998, Brain Res Rev 28:309--369{]} (https://www.sciencedirect.com/science/article/abs/pii/S0165017398000195){]} or update value representations in orbitofrontal and ventromedial prefrontal cortex {[}{[}Schultz, 2015, Neuron 86:646--664{]}(https://www.cell.com/neuron/fulltext/S0896-6273(15)00323- 1){]}. Behaviorally, this manifests as anhedonia (inability to experience pleasure from previously rewarding activities), amotivation (failure to initiate goal-directed behavior despite intact knowledge of what should be done), and a subjective sense that continued existence lacks purpose. The person in this state is abstinent by behavioral criteria but biologically incapable of deriving reward from the activities that would, in a neurotypical individual, constitute ``recovery.'' The relapse pathway under these conditions is not a matter of weak willpower or insufficient commitment to sobriety. It is a predictable consequence of living in a state of persistent reward deficit. The drug remains the only stimulus in the person's experiential repertoire capable of generating sufficient dopaminergic signal to feel worth pursuing. Abstinence demands sustained effort allocation toward a future state (long-term sobriety, reintegration into work/family, health restoration) that the person's reward circuitry cannot represent as valuable. The drug, by contrast, is an immediately available, pharmacologically reliable generator of dopamine release that temporarily restores the subjective experience of agency and positive affect. Under these conditions, relapse is not a failure of treatment --- it is evidence that treatment never addressed the underlying deficit.

\subsubsection{The Measurement Gap as Diagnostic}\label{18_restoration-the-measurement-gap-as-diagnostic}

The existing recovery literature's consistent failure to measure hedonic tone, motivational capacity, or quality of abstinence constitutes evidence that the field is not asking the right questions. A systematic review of randomized controlled trials for stimulant use disorder published between 2000 and 2020 would reveal that primary endpoints are overwhelmingly abstinence-based: percentage of drug-negative urine screens, longest period of continuous abstinence, time to first use, or proportion of participants achieving 3+ weeks of continuous abstinence. Secondary endpoints, when present, typically include retention in treatment, craving scores, and Addiction Severity Index composite scores. What is systematically absent from this literature is any validated measure of reward function. The field does not routinely assess:

\begin{itemize}
\tightlist
\item
  Hedonic capacity: the ability to experience pleasure from natural reinforcers, distinct from absence of dysphoria\\
\item
  Motivational function: the capacity to initiate and sustain goal-directed behavior in the absence of external coercion\\
\item
  Subjective sense of agency: the phenomenological experience of being the author of one's actions rather than a passive spectator\\
\item
  Quality of abstinence: whether the abstinent state is experienced as livable or as white-knuckle endurance Depression scales (Hamilton Depression Rating Scale, Beck Depression Inventory, Montgomery-Åsberg Depression Rating Scale) capture mood valence and the presence of negative affect but are not designed to measure the specific deficit in reward sensitivity and effort allocation that characterizes dopaminergic dysfunction. The Snaith-Hamilton Pleasure Scale (SHAPS) {[}\href{https://pubmed.ncbi.nlm.nih.gov/7551619/}{Snaith et al., 1995, Br J Psychiatry 167:99-- 103}{]}, though designed to assess anhedonia, focuses on consummatory pleasure rather than the anticipatory and motivational components of reward that depend on mesolimbic dopamine. No widely adopted instrument in addiction trials measures what Treadway et al.~(2012) termed ``effort-based decision-making for reward'' {[}{[}Treadway et al., 2012, J Neurosci 32:6170--6176{]} (https://www.jneurosci.org/content/32/18/6170){]} --- the willingness to expend effort for uncertain future gains, which is the operational definition of motivated behavior and the specific function that D2/D3 receptor downregulation impairs. This measurement gap is not an oversight. It reflects the field's implicit theoretical commitment to a model in which drug use is the pathology and cessation of drug use is the cure. Under this model, what happens to the person's subjective experience after drug cessation is irrelevant as long as they remain abstinent. The failure to measure reward function is thus not a methodological limitation but a revealed preference: the field has operationalized recovery as behavioral compliance (not using drugs) rather than functional restoration (being able to live a life worth living). The cost of this mismeasurement is borne by patients, who are cycled through treatment programs that achieve temporary abstinence without restoring the neurobiological capacity for sustained goal pursuit, then labeled as ``unmotivated'' or ``resistant'' when they relapse into the only state that makes them feel human.
\end{itemize}

\begin{center}\rule{0.5\linewidth}{0.5pt}\end{center}

\subsection{Section 2: What the Target Actually Is, Biologically}\label{18_restoration-section-2-what-the-target-actually-is-biologically}

\subsubsection{Mesolimbic Reward Capacity as an Integrated System}\label{18_restoration-mesolimbic-reward-capacity-as-an-integrated-system}

Mesolimbic reward capacity, properly defined, is not a single variable but an integrated system comprising at least four interdependent components: (a) D2/D3 receptor density in striatum as a proxy for postsynaptic sensitivity to dopaminergic input; (b) dopamine synthesis and release capacity in VTA/NAcc projections; (c) mitochondrial function and oxidative stress burden in dopaminergic neurons as upstream determinants of both synthesis and receptor expression; and (d) neuroinflammatory tone as a modulator of all three. These components form a cascade rather than independent variables, and intervening at only one node is likely insufficient for durable recovery.

\subsubsection{Component (a): D2/D3 Receptor Density}\label{18_restoration-component-a-d2d3-receptor-density}

D2-type receptors (D2 and D3 subtypes) in the striatum --- particularly in the nucleus accumbens and dorsal striatum --- serve as the postsynaptic transduction mechanism for dopaminergic signals encoding reward prediction error, incentive salience, and effort valuation. Binding of dopamine to D2/D3 receptors on medium spiny neurons of the indirect pathway disinhibits thalamic and cortical circuits, enabling action selection and behavioral reinforcement {[}{[}Gerfen \& Surmeier, 2011, Annu Rev Neurosci 34:441--466{]}(https://www.annualreviews.org/content/journals/10.1146/annurev-neuro- 061010-113641){]}. Critically, D2/D3 receptor availability is not a static trait but a dynamically regulated endpoint that is massively downregulated by chronic stimulant exposure. The mechanism operates across three timescales, as detailed in the companion D2 downregulation document. At the fastest timescale (hours to days), dopamine flooding triggers classical GPCR desensitization: GRK2-mediated phosphorylation, β-arrestin2 recruitment, clathrin-mediated endocytosis, and GASP-1-dependent lysosomal degradation {[}{[}Bartlett et al., 2005, PNAS 102:11521--11526{]} (https://pmc.ncbi.nlm.nih.gov/articles/PMC1183554/){]}. Unlike D1 receptors, which are efficiently recycled to the membrane {[}\href{https://pubmed.ncbi.nlm.nih.gov/9885242/}{Vickery \& von Zastrow, 1999, J Cell Biol 144:31--43}{]}, D2 and D3 receptors bind GASP- 1 with high affinity and are routed to lysosomes for irreversible destruction {[}{[}Heydorn et al., 2004, J Biol Chem 279:44989--44998{]} (https://pubmed.ncbi.nlm.nih.gov/15304494/){]}. Each episode of stimulant use thus produces a net loss of D2R protein. At intermediate timescales (weeks to months), chronic stimulant exposure induces ΔFosB, a transcription factor with an extraordinary 6--8 week half-life that persists long after the last drug exposure {[}{[}Nestler, 2008, Nat Rev Neurosci 9:844--858{]} (https://www.nature.com/articles/nrn2491){]}. ΔFosB recruits HDAC1 to the DRD2 promoter, causing histone H4 deacetylation and transcriptional silencing {[}\href{https://pubmed.ncbi.nlm.nih.gov/18632938/}{Renthal et al., 2008, J Neurosci 28:7344--7349}{]}. Simultaneously, repressive histone methylation marks (H3K27me3) accumulate at the DRD2 locus, creating a dual epigenetic lock --- deacetylation plus methylation --- that compacts chromatin and prevents compensatory upregulation of DRD2 transcription even as receptor protein is being degraded {[}\href{https://pubmed.ncbi.nlm.nih.gov/28421257/}{Montalvo-Ortiz et al., 2017, Psychopharmacology 234:2385--2398}{]}. At the slowest timescale (months to potentially permanent), DNA methylation at CpG sites in the DRD2 promoter may stabilize the repressed state. Unlike histone modifications, which turn over on timescales of hours to weeks, DNA methylation in post-mitotic neurons is extraordinarily stable and resists passive reversal {[}{[}Moore et al., 2013, Neuropsychopharmacology 38:23--38{]} (https://www.nature.com/articles/npp2012112){]}. Whether chronic stimulant use establishes persistent DRD2 promoter methylation in human striatal neurons remains an open question, but animal data and cross-sectional human studies suggest it does {[}{[}Nohesara et al., 2016, Am J Med Genet Part B 171:1180--1189{]} (https://pmc.ncbi.nlm.nih.gov/articles/PMC7115129/){]}. The PET imaging evidence is unambiguous: D2/D3 receptor downregulation in stimulant users is large (effect sizes of −0.73 to −0.81 in meta-analyses encompassing over 500 users) {[}{[}Ashok et al., 2017, JAMA Psychiatry 74:511--519{]} (https://jamanetwork.com/journals/jamapsychiatry/fullarticle/2608759){]} and shows no spontaneous recovery at 9 months of abstinence in the only longitudinal study to measure it {[}{[}Volkow et al., 2015, NeuroImage 121:20--28{]} (https://www.sciencedirect.com/science/article/abs/pii/S1053811915006461){]}. This is the single most important neurobiological finding in the addiction literature, and current treatment paradigms ignore it entirely.

\subsubsection{Component (b): Dopamine Synthesis and Release Capacity}\label{18_restoration-component-b-dopamine-synthesis-and-release-capacity}

Receptor density is meaningless if the presynaptic terminal cannot synthesize and release adequate dopamine. Chronic stimulant exposure damages the presynaptic machinery through oxidative stress, mitochondrial dysfunction, and cytoskeletal disruption. Methamphetamine in particular produces direct neurotoxicity to dopaminergic terminals through mechanisms including dopamine auto-oxidation, reactive oxygen species generation, and disruption of vesicular storage by VMAT2 inhibition {[}{[}Krasnova \& Cadet, 2009, Neuropharmacology 56 Suppl 1:83--89{]} (https://pmc.ncbi.nlm.nih.gov/articles/PMC2631950/){]}. Partial recovery of dopamine transporter density occurs at 9--17 months of abstinence {[}{[}Volkow et al., 2015, NeuroImage 121:20--28{]} (https://www.sciencedirect.com/science/article/abs/pii/S1053811915006461){]}, suggesting some degree of terminal repair. However, DAT recovery is not accompanied by normalization of dopamine release capacity: when abstinent methamphetamine users are challenged with methylphenidate (which increases synaptic dopamine by blocking reuptake), those who subsequently relapse show no increase in striatal dopamine release, whereas those who maintain abstinence show release comparable to controls {[}{[}Wang et al., 2012, Mol Psychiatry 17:918--925{]} (https://pmc.ncbi.nlm.nih.gov/articles/PMC3261322/){]}. This dissociation implies that structural repair of the terminal (DAT reappearance) does not guarantee functional restoration of dopamine synthesis and packaging capacity. Tyrosine hydroxylase, the rate-limiting enzyme in dopamine synthesis, requires tetrahydrobiopterin (BH4) as a cofactor and is subject to feedback inhibition by dopamine itself. Under conditions of chronic stimulation, TH may be downregulated, oxidatively damaged, or functionally impaired by cofactor depletion. Restoration of synthesis capacity likely requires neurotrophic support (GDNF, BDNF, CDNF) to drive transcriptional upregulation of TH, GTP cyclohydrolase I (the rate-limiting enzyme in BH4 synthesis), and VMAT2.

\subsubsection{Component (c): Mitochondrial Function and Oxidative Stress}\label{18_restoration-component-c-mitochondrial-function-and-oxidative-stress}

Mitochondrial dysfunction in dopaminergic neurons is both a consequence of stimulant exposure and an upstream driver of both dopamine synthesis deficits and D2 receptor downregulation. Methamphetamine impairs complex I of the electron transport chain {[}{[}Wernicke et al., 2010, Pharmacol Rep 62:35--53{]} (https://pubmed.ncbi.nlm.nih.gov/20360614/){]}, documenting \textasciitilde50\% reduction in striatal complex I activity in the MPP+ rat model of dopaminergic neurotoxicity; the experimental compound 9-Me-BC reversed this deficit by \textasciitilde80\% {[}{[}Wernicke et al., 2010{]} (https://pubmed.ncbi.nlm.nih.gov/20360614/){]}. Complex I impairment reduces ATP production, increases electron leak and reactive oxygen species generation, and decreases the electrochemical gradient required for mitochondrial dopamine sequestration by VMAT2 {[}{[}Lohr et al., 2017, Neurochem Res 42:2639--2662{]} (https://link.springer.com/article/10.1007/s11064-017-2259-4){]}. Oxidative stress, in turn, oxidizes dopamine to toxic quinones, damages TH and other catecholamine synthetic enzymes, and triggers inflammatory signaling cascades. The relationship to receptor downregulation is bidirectional: oxidative stress and mitochondrial dysfunction impair the energetic capacity to maintain protein synthesis, including synthesis of D2 receptors to replace those degraded via GASP-1; conversely, D2 receptor signaling itself modulates mitochondrial function and oxidative tone in striatal neurons, such that D2R loss may perpetuate mitochondrial impairment.

\subsubsection{Component (d): Neuroinflammatory Tone}\label{18_restoration-component-d-neuroinflammatory-tone}

Converging evidence links neuroinflammation to reduced dopamine synthesis and impaired reward-circuit responsivity. Inflammatory cytokines (IL-1β, IL-6, TNF-α) activate indoleamine 2,3-dioxygenase (IDO), shunting tryptophan metabolism away from serotonin synthesis toward kynurenine, which is further metabolized to quinolinic acid --- an NMDA receptor agonist and mitochondrial toxin {[}\href{https://www.nature.com/articles/nrn2297}{Dantzer et al., 2008, Nat Rev Neurosci 9:46--56}{]}. Inflammatory signaling also activates microglia and astrocytes, which can adopt neurotoxic phenotypes that impair dopaminergic neuron function {[}\href{https://www.nature.com/articles/nrn2038}{Block et al., 2007, Nat Rev Neurosci 8:57--69}{]}. Clinical trials and secondary analyses in major depression demonstrate that elevated inflammatory markers (e.g., C-reactive protein) predict differential antidepressant response, with relatively better outcomes on catecholaminergic agents (bupropion) than serotonergic agents (SSRIs) {[}{[}Uher et al., 2014, Am J Psychiatry 171:1278--1286{]} (https://psychiatryonline.org/doi/10.1176/appi.ajp.2014.14010094); \href{https://pubmed.ncbi.nlm.nih.gov/28090845/}{Jha et al., 2017, Psychol Med 47:1136--1148}{]}. In addiction specifically, Bekhbat et al.~(2022) showed that anhedonia severity in cocaine use disorder correlated with reduced functional connectivity in reward circuitry and that this deficit was mechanistically linked to inflammatory tone {[}\href{https://www.nature.com/articles/s41380-022-01715-3}{Bekhbat et al., 2022, Mol Psychiatry 27:4113--4121}{]}. The inflammatory subtype of depression/addiction may represent a distinct biological phenotype requiring dopaminergic rather than serotonergic intervention.

\subsubsection{Why the System Is Integrated}\label{18_restoration-why-the-system-is-integrated}

\subsection{The four components are not independent variables but nodes in a causal network. Oxidative stress drives mitochondrial dysfunction, which impairs dopamine synthesis, which produces compensatory D2 downregulation, which is then epigenetically locked by ΔFosB/HDAC1-mediated chromatin remodeling. Neuroinflammation amplifies oxidative stress and directly inhibits dopamine synthesis via IDO. D2 receptor loss reduces the capacity of surviving dopamine to encode reward, further demotivating the behaviors (exercise, social engagement, goal pursuit) that would otherwise stimulate endogenous neurotrophic factor release. This integration explains why single-node interventions are likely to fail. Blocking D2 degradation with a hypothetical GASP-1 inhibitor would be futile if transcription of new receptors is epigenetically silenced. Reversing histone deacetylation with an HDAC inhibitor would be insufficient if mitochondrial dysfunction prevents adequate dopamine synthesis to signal through the newly expressed receptors. Restoring dopamine synthesis with a TH upregulator would provide no benefit if postsynaptic receptor density remains at 60\% of baseline. Durable recovery of mesolimbic reward capacity therefore requires interventions that address multiple nodes simultaneously: neurotrophic support to restore synthesis capacity (exercise, GDNF-inducing agents such as ibogaine), mitochondrial restoration (complex I support, antioxidant systems), epigenetic unlocking (HDAC inhibitors), and potentially direct receptor restoration (AAV-DRD2, CRISPRa-DRD2). No current treatment protocol approaches this level of mechanistic targeting.}\label{18_restoration-the-four-components-are-not-independent-variables-but-nodes-in-a-causal-network.-oxidative-stress-drives-mitochondrial-dysfunction-which-impairs-dopamine-synthesis-which-produces-compensatory-d2-downregulation-which-is-then-epigenetically-locked-by-ux3b4fosbhdac1-mediated-chromatin-remodeling.-neuroinflammation-amplifies-oxidative-stress-and-directly-inhibits-dopamine-synthesis-via-ido.-d2-receptor-loss-reduces-the-capacity-of-surviving-dopamine-to-encode-reward-further-demotivating-the-behaviors-exercise-social-engagement-goal-pursuit-that-would-otherwise-stimulate-endogenous-neurotrophic-factor-release.-this-integration-explains-why-single-node-interventions-are-likely-to-fail.-blocking-d2-degradation-with-a-hypothetical-gasp-1-inhibitor-would-be-futile-if-transcription-of-new-receptors-is-epigenetically-silenced.-reversing-histone-deacetylation-with-an-hdac-inhibitor-would-be-insufficient-if-mitochondrial-dysfunction-prevents-adequate-dopamine-synthesis-to-signal-through-the-newly-expressed-receptors.-restoring-dopamine-synthesis-with-a-th-upregulator-would-provide-no-benefit-if-postsynaptic-receptor-density-remains-at-60-of-baseline.-durable-recovery-of-mesolimbic-reward-capacity-therefore-requires-interventions-that-address-multiple-nodes-simultaneously-neurotrophic-support-to-restore-synthesis-capacity-exercise-gdnf-inducing-agents-such-as-ibogaine-mitochondrial-restoration-complex-i-support-antioxidant-systems-epigenetic-unlocking-hdac-inhibitors-and-potentially-direct-receptor-restoration-aav-drd2-crispra-drd2.-no-current-treatment-protocol-approaches-this-level-of-mechanistic-targeting.}

\subsection{Section 3: Why Current Treatment Targets Are Actively Counterproductive}\label{18_restoration-section-3-why-current-treatment-targets-are-actively-counterproductive}

\subsubsection{Receptor Blockade: Deepening the Anhedonic Trap}\label{18_restoration-receptor-blockade-deepening-the-anhedonic-trap}

Receptor blockade strategies --- exemplified by naltrexone for opioid use disorder and the failed antipsychotic trials in cocaine use disorder --- operate on the logic that if drug reward is mediated through a receptor, blocking that receptor will eliminate the incentive to use. For opioids, where the therapeutic target (mu-opioid receptor) is distinct from the reward-circuit substrate (dopaminergic signaling), this approach has shown modest efficacy. For stimulants, where the drug directly potentiates dopamine signaling and the pathology is D2/D3 receptor downregulation, receptor blockade is mechanistically absurd and clinically harmful. Antipsychotics block D2 receptors with high affinity, acutely reducing the dopaminergic signal generated by both cocaine and natural rewards. In the short term, this might suppress drug-induced euphoria; in the medium to long term, it worsens the underlying deficit. Chronic D2 blockade triggers compensatory upregulation of D2 receptors (the mechanism underlying antipsychotic tolerance and dopamine supersensitivity psychosis), but this upregulation occurs in the context of continued pharmacological antagonism and reverses when the drug is discontinued. More importantly, D2 blockade impairs the encoding of natural rewards --- food becomes less pleasant, social interaction less rewarding, goal accomplishment less satisfying --- deepening the anhedonic state that drove self-medication in the first place. A meta-analysis of 12 randomized controlled trials of antipsychotics for cocaine use disorder found no reduction in cocaine use and high dropout rates {[}{[}Indave et al., 2016, Cochrane Database Syst Rev, CD006306{]} (https://www.cochranelibrary.com/cdsr/doi/10.1002/14651858.CD006306.pub3/full){]}. The predictable outcome of blocking the only functional dopaminergic signaling remaining in a person with severe D2/D3 downregulation is not abstinence but despair.

\subsubsection{Abstinence-Only Approaches: Return to the Pathological Baseline}\label{18_restoration-abstinence-only-approaches-return-to-the-pathological-baseline}

Abstinence-only interventions that do not address the underlying dopaminergic deficit produce a state biologically indistinguishable from the pre-drug anhedonic baseline that drove self-medication in the first place. This is most clearly illustrated in patients who began stimulant use as self-medication for undiagnosed ADHD or treatment-resistant anhedonic depression --- a clinical trajectory detailed in the companion TRAAD document {[}{[}Anonymous, 2025, The TRAAD Spectrum: A Single-Case Analysis of Dopamine-Responsive Depression Beyond Categorical Diagnosis{]} (placeholder\_reference){]}. For these individuals, the decade preceding stimulant use was characterized by profound anhedonia, executive paralysis, inability to initiate goal-directed behavior, and pervasive suicidal ideation maintained only by ethical reasoning rather than any positive motivation to continue living. Stimulant exposure provided the first experience of functional dopaminergic tone: work became possible, social interaction became tolerable, the future became representable as something other than an extension of present misery. Tolerance and escalation followed, culminating in dependence and medical crisis. Abstinence-enforcement in this population does not restore a healthy baseline --- it returns the person to the pathological pre-drug state, now compounded by neuroadaptive changes (D2/D3 downregulation, potential receptor supersensitivity in some circuits, conditioned associations between drug cues and relief) that may make the anhedonic state even more severe than it was initially. The clinical literature terms this ``protracted withdrawal'' and attributes it to neuroadaptive consequences of chronic drug exposure. Mechanistically, it is the reemergence of the original dopamine-deficit state, no longer masked by exogenous dopamine augmentation, and now worsened by stimulant-induced D2 receptor loss. Demanding indefinite abstinence from a person in this state, without providing any intervention to restore reward capacity, is not treatment. It is a demand that they endure a neurobiological condition (severe anhedonia, amotivation, inability to encode the value of future goals) indefinitely, with no pharmacological support and often with the explicit message that any return to dopaminergic augmentation --- even under medical supervision --- constitutes moral failure. The predictable result is relapse or suicide.

\subsubsection{Coercive Approaches: Stress-Mediated Dopaminergic Damage}\label{18_restoration-coercive-approaches-stress-mediated-dopaminergic-damage}

\subsection{Coercive treatment approaches --- mandated residential programs, court-ordered abstinence monitoring, therapeutic communities employing confrontation and public humiliation, or family interventions involving threats of abandonment --- generate chronic psychosocial stress. The neurobiological consequences of chronic stress are well-characterized and directly counterproductive to dopaminergic restoration. Prefrontal dopamine signaling follows an inverted-U relationship with arousal and stress: moderate levels optimize working memory and cognitive control, while excessive catecholamine release (driven by sustained activation of the locus coeruleus and hypothalamic-pituitary-adrenal axis) impairs prefrontal function and shifts control to more primitive subcortical systems {[}Arnsten, 2009, Biol Psychiatry{]}. Chronic stress elevates cortisol, which in rodent models reduces dopamine synthesis in mesocortical projections, impairs D1 receptor signaling in prefrontal cortex, and exacerbates the very executive dysfunction (impaired planning, inability to delay gratification, perseverative errors) that is labeled as evidence of ``lack of commitment to recovery.'' Neuroinflammation is also stress-sensitive. Chronic psychosocial stress activates microglia, increases pro-inflammatory cytokine expression in hippocampus and prefrontal cortex, and elevates peripheral markers of inflammation (CRP, IL-6) that have been shown to predict poor response to serotonergic antidepressants and preferential response to dopaminergic agents {[}Bekhbat et al., 2022; Uher et al., 2014{]}. Inflammatory activation, in turn, impairs dopamine synthesis via IDO-mediated tryptophan shunting and quinolinic acid accumulation, creating a vicious cycle in which stress-induced inflammation worsens dopaminergic deficit, which worsens anhedonia and amotivation, which is interpreted as willful non-compliance, which triggers more coercive intervention and more stress. The most punitive interventions --- public shaming, forced gratitude exercises, indefinite confinement contingent on behavioral compliance --- are thus not merely ethically problematic but neurobiologically damaging. By sustaining HPA axis activation and neuroinflammatory tone, they actively deepen the dopaminergic deficit that underlies the disorder they purport to treat. The paradox is complete: the interventions designed to compel abstinence through aversive consequences are the interventions most likely to worsen the neurobiological substrate that makes abstinence unlivable.}\label{18_restoration-coercive-treatment-approaches-mandated-residential-programs-court-ordered-abstinence-monitoring-therapeutic-communities-employing-confrontation-and-public-humiliation-or-family-interventions-involving-threats-of-abandonment-generate-chronic-psychosocial-stress.-the-neurobiological-consequences-of-chronic-stress-are-well-characterized-and-directly-counterproductive-to-dopaminergic-restoration.-prefrontal-dopamine-signaling-follows-an-inverted-u-relationship-with-arousal-and-stress-moderate-levels-optimize-working-memory-and-cognitive-control-while-excessive-catecholamine-release-driven-by-sustained-activation-of-the-locus-coeruleus-and-hypothalamic-pituitary-adrenal-axis-impairs-prefrontal-function-and-shifts-control-to-more-primitive-subcortical-systems-arnsten-2009-biol-psychiatry.-chronic-stress-elevates-cortisol-which-in-rodent-models-reduces-dopamine-synthesis-in-mesocortical-projections-impairs-d1-receptor-signaling-in-prefrontal-cortex-and-exacerbates-the-very-executive-dysfunction-impaired-planning-inability-to-delay-gratification-perseverative-errors-that-is-labeled-as-evidence-of-lack-of-commitment-to-recovery.-neuroinflammation-is-also-stress-sensitive.-chronic-psychosocial-stress-activates-microglia-increases-pro-inflammatory-cytokine-expression-in-hippocampus-and-prefrontal-cortex-and-elevates-peripheral-markers-of-inflammation-crp-il-6-that-have-been-shown-to-predict-poor-response-to-serotonergic-antidepressants-and-preferential-response-to-dopaminergic-agents-bekhbat-et-al.-2022-uher-et-al.-2014.-inflammatory-activation-in-turn-impairs-dopamine-synthesis-via-ido-mediated-tryptophan-shunting-and-quinolinic-acid-accumulation-creating-a-vicious-cycle-in-which-stress-induced-inflammation-worsens-dopaminergic-deficit-which-worsens-anhedonia-and-amotivation-which-is-interpreted-as-willful-non-compliance-which-triggers-more-coercive-intervention-and-more-stress.-the-most-punitive-interventions-public-shaming-forced-gratitude-exercises-indefinite-confinement-contingent-on-behavioral-compliance-are-thus-not-merely-ethically-problematic-but-neurobiologically-damaging.-by-sustaining-hpa-axis-activation-and-neuroinflammatory-tone-they-actively-deepen-the-dopaminergic-deficit-that-underlies-the-disorder-they-purport-to-treat.-the-paradox-is-complete-the-interventions-designed-to-compel-abstinence-through-aversive-consequences-are-the-interventions-most-likely-to-worsen-the-neurobiological-substrate-that-makes-abstinence-unlivable.}

\subsection{Section 4: What Treatment Aimed at the Correct Target Would Look Like}\label{18_restoration-section-4-what-treatment-aimed-at-the-correct-target-would-look-like}

This section outlines principles, not protocols. A treatment framework oriented toward restoration of mesolimbic reward capacity would differ from current practice in four fundamental domains: outcome measurement, pharmacological evaluation criteria, intervention portfolio, and maintenance philosophy.

\subsubsection{Principle (a): Primary Outcome Measures Must Assess Reward Function}\label{18_restoration-principle-a-primary-outcome-measures-must-assess-reward-function}

The primary outcome measure in any trial claiming to treat stimulant use disorder should be a validated index of hedonic capacity and motivational function, not days abstinent. Abstinence can remain a secondary endpoint --- it is not irrelevant --- but it cannot be the definition of success if success is defined as restoration of the capacity to live a life worth living. The measurement instrument does not yet exist in standardized, widely adopted form. What is needed is a brief, validated scale that captures:

\begin{enumerate}
\def\labelenumi{\arabic{enumi}.}
\tightlist
\item
  Consummatory pleasure: the ability to experience pleasure from natural rewards (food, social contact, aesthetic experiences) in the moment\\
\item
  Anticipatory pleasure: the ability to prospectively simulate future positive events and experience motivation to pursue them\\
\item
  Effort-based decision-making: willingness to expend effort for uncertain future gains, quantifiable through tasks like the Effort Expenditure for Rewards Task (EEfRT) {[}Treadway et al., 2012{]}\\
\item
  Subjective agency: the phenomenological sense of being an active author of one's behavior rather than a passive spectator to one's own failure\\
\item
  Goal pursuit: the ability to formulate, maintain, and work toward medium-term goals (weeks to months) despite setbacks Existing instruments approach subsets of these constructs. The Temporal Experience of Pleasure Scale (TEPS) distinguishes anticipatory from consummatory pleasure. The EEfRT provides a behavioral assay of effort allocation under reward uncertainty. The Apathy Evaluation Scale and Dimensional Apathy Scale assess motivational deficits. None integrate all five domains in a format suitable for repeated administration as a trial endpoint. Development and validation of such an instrument should be a priority comparable to the development of standardized depression and anxiety scales. Until the field can reliably measure what it claims to be restoring, treatment evaluation will remain a measurement of compliance (not using drugs, attending sessions) rather than an assessment of function.
\end{enumerate}

\subsubsection{Principle (b): Pharmacological Interventions Must Be Evaluated for Effect on}\label{18_restoration-principle-b-pharmacological-interventions-must-be-evaluated-for-effect-on}

Reward Circuitry Pharmacological interventions should be evaluated for their effect on D2/D3 receptor availability by PET, on dopamine synthesis capacity (as inferred from FDOPA PET or CSF metabolite profiles), and on mitochondrial/oxidative markers (systemic surrogates: lactate, oxidized glutathione; research settings: ³¹P magnetic resonance spectroscopy for ATP/PCr ratios), not merely on self-reported craving or urine drug screens. This standard would transform the evidence base. Naltrexone, widely prescribed despite minimal efficacy for stimulant use disorder, has never been evaluated for its effect on striatal D2/D3 receptor availability in stimulant users. Bupropion, which shows modest benefit in meta-analyses, has not been assessed for whether it reverses or attenuates stimulant-induced D2 downregulation. The only intervention with human PET evidence for increasing D2/D3 receptor availability is structured aerobic exercise {[}Robertson et al., 2016, Neuropsychopharmacology: 13.89\% increase in striatal D2/D3 BPND after 8 weeks, 1 hour/session, 3×/week{]}. For novel interventions --- HDAC inhibitors, GDNF-inducing agents (ibogaine), mitochondrial cofactor supplementation, or eventually gene therapy --- the evidentiary standard should be: does this increase D2/D3 receptor availability, does it increase dopamine release capacity in response to natural rewards, does it restore the functional connectivity of mesolimbic circuits (NAcc-VTA, NAcc-vmPFC, striatum-dlPFC) that is disrupted in stimulant users? If the answer is no, the intervention may suppress drug use through aversive conditioning or external contingencies, but it is not addressing the deficit.

\subsubsection{Principle (c): Intervention Portfolios Must Address the Integrated System}\label{18_restoration-principle-c-intervention-portfolios-must-address-the-integrated-system}

Because mesolimbic reward capacity is an integrated system (D2 receptor density, dopamine synthesis, mitochondrial function, inflammatory tone), durable recovery likely requires interventions addressing all four components simultaneously rather than sequentially. \textbf{Neurotrophic support:} Exercise is the only intervention with human evidence for D2 receptor upregulation and should be a mandatory, structured, supervised component of every stimulant use disorder treatment program. One hour of aerobic exercise, three times per week, for a minimum of 8 weeks is the evidence-based threshold {[}{[}Robertson et al., 2016, Neuropsychopharmacology 41:1629--1636{]} (https://pmc.ncbi.nlm.nih.gov/articles/PMC4832026/), reporting 13.89\% increase in striatal D2/D3 BPND measured by {[}¹F{]}fallypride PET{]}. This is not adjunctive therapy or a wellness add-on --- it is the only demonstrated method for reversing the core pathology. GDNF-inducing agents, particularly ibogaine, represent a pharmacological route to neurotrophic restoration. The mechanism is well-characterized in rodents (VTA GDNF mRNA upregulation, autoregulatory positive-feedback loop persisting weeks beyond drug clearance) {[}{[}He et al., 2005, J Neurosci 25:619--628{]} (https://www.jneurosci.org/content/25/3/619); \href{https://faseb.onlinelibrary.wiley.com/doi/10.1096/fj.06-6394fje}{He \& Ron, 2006, FASEB J 20:E1820-- E1827}{]}, and open- label human data from Brazil show 61\% abstinence at follow-up with median durations of 5.5--8.4 months {[}{[}Schenberg et al., 2014, J Psychopharmacol 28:993--1000{]} (https://pubmed.ncbi.nlm.nih.gov/25271214/){]}. No PET imaging of D2/D3 receptors before and after ibogaine has been conducted in humans, which is the most consequential gap in the literature. A {[}¹¹C{]}raclopride or {[}¹F{]}fallypride study in chronic stimulant users pre- and post-ibogaine, conducted under contemporary cardiac risk-management protocols (CYP2D6 genotyping, magnesium loading, telemetry), would directly test whether the clinical improvements are accompanied by receptor restoration. \textbf{Mitochondrial restoration:} Preclinical evidence suggests coenzyme Q10, N- acetylcysteine (antioxidant, glutathione precursor), and potentially 9-Me-BC (complex I upregulator, TH inducer, though entirely preclinical with no human safety data) could address mitochondrial and oxidative components {[}\href{https://pubmed.ncbi.nlm.nih.gov/20360614/}{Wernicke et al., 2010, Pharmacol Rep 62:35--53}{]}. These interventions are safe, inexpensive, and mechanistically plausible but have not been tested in adequately powered trials with D2/D3 PET endpoints. \textbf{Epigenetic unlocking:} HDAC inhibitors reverse ΔFosB/HDAC1-mediated transcriptional repression of DRD2 in rodent models {[}{[}McClarty et al., 2021, Front Neurosci 15:674745{]} (https://www.frontiersin.org/journals/neuroscience/articles/10.3389/fnins.2021.674745 /full), reporting CI-994 increased H3K27ac at Drd2 promoter, restored D2R expression and function{]}. Clinical development of brain-penetrant, selective HDAC inhibitors for addiction is feasible within 8--12 years using compounds already in preclinical development (ACY-738, largazole, PB94). No such trial is currently funded or planned. \textbf{Direct receptor restoration (future):} AAV-DRD2 gene therapy and CRISPRa-mediated DRD2 upregulation are technically feasible and have proof-of-concept in rodent addiction models {[}{[}Thanos et al., 2001, Gene Ther 8:1479--1486{]} (https://www.nature.com/articles/3301537); {[}Thanos et al., 2004, Synapse 54:119--128{]} (https://pubmed.ncbi.nlm.nih.gov/15352133/), AAV-DRD2 reduced alcohol self- administration by 43--64\%{]}. The Upstaza precedent (EMA-approved AAV-AADC for aromatic L-amino acid decarboxylase deficiency, administered via direct intrastriatal injection) {[}{[}Hoy, 2022, Mol Diagn Ther 26:603--610{]} (https://link.springer.com/article/10.1007/s40291-022-00608-7){]} demonstrates that stereotactic AAV delivery to human striatum is viable. Blood-brain barrier technologies (FUS-BBBO, engineered AAV capsids, brain-targeted LNPs) are advancing rapidly enough that systemic delivery of receptor-restorative gene therapy could be feasible within 15 years. This timeline is realistic only if institutional funding is redirected toward addiction, which currently generates negligible pharmaceutical revenue compared to oncology and neurodegeneration.

\subsubsection{Principle (d): Maintenance Rather Than Abstinence Should Be the Default Frame for}\label{18_restoration-principle-d-maintenance-rather-than-abstinence-should-be-the-default-frame-for}

\subsection{Severe Dopaminergic Deficit For individuals with severe, treatment-refractory dopaminergic deficits --- those who meet criteria for the TRAAD phenotype (treatment-resistant anhedonic and amotivational depression), those with baseline D2/D3 receptor availability in the lowest quartile, those with documented complete non-response to serotonergic agents and robust response to dopaminergic agents --- maintenance pharmacotherapy should be the default frame, not abstinence. The analogy to insulin in type 1 diabetes and levodopa in Parkinson's disease is precise, not rhetorical. In type 1 diabetes, pancreatic beta-cell loss creates an absolute deficiency of insulin; withholding insulin does not motivate the pancreas to recover, it produces hyperglycemia, ketoacidosis, and death. In Parkinson's disease, nigrostriatal dopaminergic neuron loss creates a motor and cognitive deficit; withholding levodopa does not stimulate neuronal regeneration, it produces rigidity, akinesia, and progressive disability. In both conditions, maintenance pharmacotherapy is recognized as physiological replacement rather than dependence, and the idea of withholding the effective agent until the patient demonstrates sufficient moral worthiness is immediately recognizable as monstrous. In severe dopaminergic-deficit depression or ADHD with stimulant response, the logic is identical. If the person cannot initiate goal-directed behavior, cannot experience reward from natural reinforcers, cannot maintain motivation over days to weeks without dopaminergic augmentation, and has exhausted all non-stimulant pharmacological options (bupropion, modafinil, atomoxetine, MAOIs in appropriate cases) and non-pharmacological interventions (structured exercise, intensive behavioral activation), then supervised, dose-controlled maintenance stimulant therapy is not ``enabling addiction'' --- it is providing the minimum neurochemical substrate required for agency. The current paradigm pathologizes this as drug-seeking and labels the physiological need for dopaminergic tone as evidence of moral deficiency. The correct framing is that some fraction of people have dopaminergic deficits severe enough that no currently available non-stimulant intervention can restore reward function to a livable threshold, and that for these individuals, chronic pharmacological support is a legitimate medical treatment, not a moral failure. Diabetes and Parkinson's frameworks have accepted this for insulin and levodopa; addiction medicine must accept it for dopaminergic agents when they are the only effective intervention for a documented, severe, treatment-refractory neurobiological deficit. This does not mean unrestricted access or dismissal of escalation risk. It means structured, monitored maintenance within a harm-reduction framework: dose caps, prescription monitoring, contingent intensification of behavioral support if dose creep occurs, but not blanket prohibition based on the principle that needing the drug disqualifies one from receiving it.}\label{18_restoration-severe-dopaminergic-deficit-for-individuals-with-severe-treatment-refractory-dopaminergic-deficits-those-who-meet-criteria-for-the-traad-phenotype-treatment-resistant-anhedonic-and-amotivational-depression-those-with-baseline-d2d3-receptor-availability-in-the-lowest-quartile-those-with-documented-complete-non-response-to-serotonergic-agents-and-robust-response-to-dopaminergic-agents-maintenance-pharmacotherapy-should-be-the-default-frame-not-abstinence.-the-analogy-to-insulin-in-type-1-diabetes-and-levodopa-in-parkinsons-disease-is-precise-not-rhetorical.-in-type-1-diabetes-pancreatic-beta-cell-loss-creates-an-absolute-deficiency-of-insulin-withholding-insulin-does-not-motivate-the-pancreas-to-recover-it-produces-hyperglycemia-ketoacidosis-and-death.-in-parkinsons-disease-nigrostriatal-dopaminergic-neuron-loss-creates-a-motor-and-cognitive-deficit-withholding-levodopa-does-not-stimulate-neuronal-regeneration-it-produces-rigidity-akinesia-and-progressive-disability.-in-both-conditions-maintenance-pharmacotherapy-is-recognized-as-physiological-replacement-rather-than-dependence-and-the-idea-of-withholding-the-effective-agent-until-the-patient-demonstrates-sufficient-moral-worthiness-is-immediately-recognizable-as-monstrous.-in-severe-dopaminergic-deficit-depression-or-adhd-with-stimulant-response-the-logic-is-identical.-if-the-person-cannot-initiate-goal-directed-behavior-cannot-experience-reward-from-natural-reinforcers-cannot-maintain-motivation-over-days-to-weeks-without-dopaminergic-augmentation-and-has-exhausted-all-non-stimulant-pharmacological-options-bupropion-modafinil-atomoxetine-maois-in-appropriate-cases-and-non-pharmacological-interventions-structured-exercise-intensive-behavioral-activation-then-supervised-dose-controlled-maintenance-stimulant-therapy-is-not-enabling-addiction-it-is-providing-the-minimum-neurochemical-substrate-required-for-agency.-the-current-paradigm-pathologizes-this-as-drug-seeking-and-labels-the-physiological-need-for-dopaminergic-tone-as-evidence-of-moral-deficiency.-the-correct-framing-is-that-some-fraction-of-people-have-dopaminergic-deficits-severe-enough-that-no-currently-available-non-stimulant-intervention-can-restore-reward-function-to-a-livable-threshold-and-that-for-these-individuals-chronic-pharmacological-support-is-a-legitimate-medical-treatment-not-a-moral-failure.-diabetes-and-parkinsons-frameworks-have-accepted-this-for-insulin-and-levodopa-addiction-medicine-must-accept-it-for-dopaminergic-agents-when-they-are-the-only-effective-intervention-for-a-documented-severe-treatment-refractory-neurobiological-deficit.-this-does-not-mean-unrestricted-access-or-dismissal-of-escalation-risk.-it-means-structured-monitored-maintenance-within-a-harm-reduction-framework-dose-caps-prescription-monitoring-contingent-intensification-of-behavioral-support-if-dose-creep-occurs-but-not-blanket-prohibition-based-on-the-principle-that-needing-the-drug-disqualifies-one-from-receiving-it.}

\subsection{Section 5: The Measurement Gap and What Closing It Requires}\label{18_restoration-section-5-the-measurement-gap-and-what-closing-it-requires}

The transition from the current paradigm to a reward-restoration framework requires building what does not yet exist. Four specific gaps are identifiable and addressable:

\subsubsection{(a) A Validated Clinical Instrument for Hedonic and Motivational Capacity}\label{18_restoration-a-a-validated-clinical-instrument-for-hedonic-and-motivational-capacity}

The field needs a brief (≤20 items), validated, psychometrically sound scale that can serve as a primary endpoint in addiction trials and that captures the specific agency-and-reward deficit described throughout this document. This instrument must be distinct from general depression scales, which assess mood valence and negative affect but do not reliably measure the motivational and anticipatory reward components that depend on dopaminergic signaling. Candidate constructs to integrate: consummatory pleasure (TEPS consummatory subscale), anticipatory pleasure (TEPS anticipatory subscale) {[}\href{https://pubmed.ncbi.nlm.nih.gov/17100538/}{Gard et al., 2006, J Abnorm Psychol 115:817--825}{]}, effort-based decision-making (computational modeling of EEfRT performance or similar effort- discounting tasks) {[}{[}Treadway et al., 2012, J Neurosci 32:6170--6176{]} (https://www.jneurosci.org/content/32/18/6170){]}, subjective sense of agency (phenomenological interview items), and goal pursuit capacity (self-report of ability to formulate and work toward medium-term goals). The scale should show sensitivity to dopaminergic interventions (stimulants, bupropion, pramipexole) and insensitivity to serotonergic interventions in the absence of comorbid anxiety/rumination, and it should predict relapse better than craving scales or abstinence duration. Development would require: item generation from patient interviews and mechanistic theory, psychometric validation in stimulant-using samples, test-retest reliability assessment, convergent validity with behavioral tasks (EEfRT) and divergent validity from depression/anxiety scales, and predictive validity for relapse and functional outcomes. Estimated timeline: 3--5 years from initial development to widespread adoption, if adequately funded.

\subsubsection{(b) Longitudinal D2/D3 PET Studies Beyond 9 Months}\label{18_restoration-b-longitudinal-d2d3-pet-studies-beyond-9-months}

No human study has followed D2/D3 receptors longitudinally beyond 9 months, and the single study at 9 months found no recovery whatsoever. Whether spontaneous recovery occurs at 18 months, 2 years, or 5 years is entirely unknown. Whether exercise- induced increases persist after cessation of the exercise protocol is unknown. Whether ibogaine, HDAC inhibitors, or mitochondrial cofactors produce measurable D2/D3 increases in humans is unknown because the experiments have not been done. A comprehensive research agenda would include:

\begin{enumerate}
\def\labelenumi{\arabic{enumi}.}
\tightlist
\item
  \textbf{Natural history study:} {[}¹¹C{]}raclopride or {[}¹F{]}fallypride PET in chronic stimulant users at baseline, 9 months, 18 months, 36 months, and 60 months of continuous abstinence (n ≥ 50 per timepoint, with retention strategies to minimize dropout). This establishes whether, when, and to what extent spontaneous recovery occurs.\\
\item
  \textbf{Exercise intervention with extended follow-up:} Replication of Robertson et al.~(2016) with a larger sample (n ≥ 60), active control (health education), and PET imaging at baseline, end of 8-week intervention, and 6-month post-intervention follow-up to assess durability.\\
\item
  \textbf{Ibogaine PET study:} Single-dose ibogaine (under contemporary cardiac risk- management protocols) with {[}¹F{]}fallypride PET at baseline, 1 month, 3 months, and 6 months post-treatment in cocaine or methamphetamine users (n ≥ 30). This is the single most consequential missing experiment in the ibogaine literature.\\
\item
  \textbf{HDAC inhibitor proof-of-concept:} If and when a brain-penetrant selective HDAC inhibitor enters Phase 1 testing, a Phase 2a study in stimulant users with D2/D3 PET as the primary endpoint (baseline and post-treatment) would provide direct mechanistic validation. None of these studies are currently funded. The barrier is not technical --- PET imaging with {[}¹¹C{]}raclopride and {[}¹F{]}fallypride is routine --- but institutional: addiction research budgets are dominated by behavioral interventions and medication repurposing trials, with minimal allocation to mechanistic neuroimaging or intervention development targeting the receptor deficit itself.
\end{enumerate}

\subsubsection{(c) Integrated Biomarker Panels Combining PET, Mitochondrial Function, and}\label{18_restoration-c-integrated-biomarker-panels-combining-pet-mitochondrial-function-and}

Neuroinflammation PET imaging of D2/D3 receptors provides a single endpoint --- postsynaptic receptor availability --- but does not distinguish whether deficits are driven by receptor degradation, transcriptional repression, presynaptic synthesis failure, mitochondrial dysfunction, or inflammatory tone. Comprehensive assessment of the integrated system requires multimodal biomarker panels:

\begin{itemize}
\tightlist
\item
  \textbf{PET:} {[}¹¹C{]}raclopride or {[}¹F{]}fallypride for D2/D3 receptors; {[}¹F{]}DOPA for dopamine synthesis capacity; {[}¹¹C{]}PE2I or similar for DAT\\
\item
  \textbf{Peripheral markers of inflammation:} high-sensitivity CRP, IL-6, TNF-α, kynurenine/tryptophan ratio (as a surrogate for IDO activity)\\
\item
  \textbf{Mitochondrial function surrogates:} lactate (plasma or CSF), oxidized/reduced glutathione ratio, urinary 8-oxo-dG (DNA oxidation marker); in research settings, ³¹P-MRS for brain ATP and phosphocreatine\\
\item
  \textbf{Epigenetic markers:} This is the most challenging. DRD2 promoter methylation and histone modification status require brain tissue or at minimum CSF, neither of which is feasible in large samples. Peripheral blood DNA methylation at the DRD2 locus has been used as a surrogate, but the correlation between peripheral and central methylation is imperfect. The goal is to stratify patients by mechanistic subtype: high inflammation/low D2 (inflammatory subtype, may respond preferentially to anti-inflammatory interventions or ibogaine), low inflammation/low D2/low FDOPA (presynaptic synthesis deficit, may benefit from mitochondrial support or TH upregulation), low D2/normal FDOPA (pure receptor deficit, candidate for HDAC inhibitor or gene therapy). Precision medicine in addiction requires precision measurement of the deficit.
\end{itemize}

\subsubsection{(d) Animal Models Defining Recovery as Restoration of Natural Reward-Seeking}\label{18_restoration-d-animal-models-defining-recovery-as-restoration-of-natural-reward-seeking}

\subsection{The vast majority of rodent addiction models define ``recovery'' as suppression of drug self-administration or extinction of drug-seeking behavior. The dependent variable is lever presses for cocaine, alcohol, or heroin, and interventions are evaluated for their ability to reduce this operant response. What is systematically not measured is whether the animal, having stopped pressing the drug lever, shows restoration of natural reward-seeking behavior: foraging for palatable food, exploration of novel environments, social interaction, wheel running (in species that spontaneously run). A recovery-oriented animal model would assess both suppression of drug-seeking and restoration of natural reward engagement. For example: after chronic cocaine or methamphetamine exposure, do animals show reduced sucrose preference, reduced social interaction time, reduced novel object exploration, reduced voluntary wheel running? After a candidate intervention (exercise, GDNF administration, HDAC inhibitor, AAV- DRD2), does natural reward-seeking recover to baseline levels, and does this correlate with striatal D2/D3 receptor density as measured by ex vivo autoradiography or in vivo PET? The distinction matters because an intervention could suppress drug-seeking through aversive conditioning, receptor blockade, or generalized behavioral suppression without restoring the capacity for natural reward. Such an intervention would be scored as successful in current animal models (drug lever presses reduced) but would fail the test of human functional recovery (can the person derive pleasure from food, work, relationships?). Reorienting animal models toward dual endpoints --- suppression of drug-seeking plus restoration of natural reward-seeking --- would produce a preclinical evidence base more aligned with the clinical goal.}\label{18_restoration-the-vast-majority-of-rodent-addiction-models-define-recovery-as-suppression-of-drug-self-administration-or-extinction-of-drug-seeking-behavior.-the-dependent-variable-is-lever-presses-for-cocaine-alcohol-or-heroin-and-interventions-are-evaluated-for-their-ability-to-reduce-this-operant-response.-what-is-systematically-not-measured-is-whether-the-animal-having-stopped-pressing-the-drug-lever-shows-restoration-of-natural-reward-seeking-behavior-foraging-for-palatable-food-exploration-of-novel-environments-social-interaction-wheel-running-in-species-that-spontaneously-run.-a-recovery-oriented-animal-model-would-assess-both-suppression-of-drug-seeking-and-restoration-of-natural-reward-engagement.-for-example-after-chronic-cocaine-or-methamphetamine-exposure-do-animals-show-reduced-sucrose-preference-reduced-social-interaction-time-reduced-novel-object-exploration-reduced-voluntary-wheel-running-after-a-candidate-intervention-exercise-gdnf-administration-hdac-inhibitor-aav--drd2-does-natural-reward-seeking-recover-to-baseline-levels-and-does-this-correlate-with-striatal-d2d3-receptor-density-as-measured-by-ex-vivo-autoradiography-or-in-vivo-pet-the-distinction-matters-because-an-intervention-could-suppress-drug-seeking-through-aversive-conditioning-receptor-blockade-or-generalized-behavioral-suppression-without-restoring-the-capacity-for-natural-reward.-such-an-intervention-would-be-scored-as-successful-in-current-animal-models-drug-lever-presses-reduced-but-would-fail-the-test-of-human-functional-recovery-can-the-person-derive-pleasure-from-food-work-relationships.-reorienting-animal-models-toward-dual-endpoints-suppression-of-drug-seeking-plus-restoration-of-natural-reward-seeking-would-produce-a-preclinical-evidence-base-more-aligned-with-the-clinical-goal.}

\subsection{Conclusion: The Logical Necessity of Paradigm Shift}\label{18_restoration-conclusion-the-logical-necessity-of-paradigm-shift}

The argument presented here is constructed such that a clinician or researcher who accepts the neurobiology cannot dismiss the treatment implications without logical inconsistency. The premises are:

\begin{enumerate}
\def\labelenumi{\arabic{enumi}.}
\tightlist
\item
  Chronic stimulant use produces massive, persistent D2/D3 receptor downregulation through GASP-1 degradation, ΔFosB/HDAC1 transcriptional repression, and epigenetic silencing.\\
\item
  D2/D3 receptor downregulation impairs the capacity of natural rewards to generate sufficient dopaminergic signal to sustain motivated behavior.\\
\item
  No spontaneous recovery of D2/D3 receptors has been documented in humans at 9 months of abstinence.\\
\item
  Abstinence without reward restoration produces an anhedonic state in which relapse is structurally inevitable rather than a moral failure.\\
\item
  Current treatment paradigms do not measure reward function, do not target the receptor deficit, and in some cases (receptor blockade, coercive stress-inducing interventions) actively worsen the underlying pathology. If these premises are accepted --- and they are not speculative but empirically established through PET imaging, molecular neuroscience, and longitudinal outcome studies --- then the current paradigm's endpoint choices (abstinence, craving suppression) are revealed as what they are: arbitrary conventions that happen to be measurable, mistaken for the thing that actually matters. The thing that actually matters is whether the person, at the end of treatment, can derive sufficient reward from natural reinforcers to make continued existence feel worthwhile. If the answer is no --- if they are abstinent but anhedonic, drug-free but unable to initiate goal-directed behavior, compliant but experiencing life as a joyless endurance test --- then treatment has not succeeded. It has produced behavioral compliance without functional restoration. The field's failure to measure this distinction, to target the neurobiological substrate that determines it, and to develop interventions that address the integrated system of dopaminergic dysfunction is not a minor methodological oversight. It is a conceptual error that has shaped decades of ineffective treatment, condemned millions of people to cycles of relapse and re-incarceration, and systematically misdirected research funding away from the interventions (exercise programs, GDNF-inducing agents, HDAC inhibitors, gene therapy, maintenance pharmacotherapy for severe deficits) that mechanistic understanding would predict to be effective. The path forward requires building what does not yet exist: validated instruments for hedonic and motivational capacity, longitudinal PET studies extending beyond 9 months, integrated biomarker panels, animal models that define recovery as restoration of natural reward rather than suppression of drug-seeking, and a clinical culture that recognizes restoration of the capacity to live a life worth living as the legitimate goal of treatment. These are not aspirational ideals --- they are the minimum requirements for a scientifically coherent approach to a disorder whose core pathology is now well-characterized and whose effective treatment requires addressing that pathology directly rather than managing its behavioral symptoms. The technologies exist. The neuroscience exists. What is missing is the institutional commitment to redirect them toward addiction --- a condition that affects millions but generates less pharmaceutical revenue than the diseases these platforms were built to treat. This is not a failure of science. It is a failure of prioritization, and it is a failure we have the knowledge and tools to correct.
\end{enumerate}

\begin{center}\rule{0.5\linewidth}{0.5pt}\end{center}

\subsection{References}\label{18_restoration-references}

Arnsten AFT (2009). The emerging neurobiology of attention-deficit/hyperactivity disorder: The key role of the prefrontal association cortex. \emph{Biological Psychiatry} 65:1168--1176. \url{https://pmc.ncbi.nlm.nih.gov/articles/PMC2894421/} Ashok AH, Mizuno Y, Volkow ND, Howes OD (2017). Association of stimulant use with dopaminergic alterations in users of cocaine, amphetamine, or methamphetamine: A systematic review and meta-analysis. \emph{JAMA Psychiatry} 74:511--519. \url{https://jamanetwork.com/journals/jamapsychiatry/fullarticle/2608759} Bartlett SE, Enquist J, Hopf FW, et al.~(2005). Dopamine responsiveness is regulated by targeted sorting of D2 receptors. \emph{Proceedings of the National Academy of Sciences} 102:11521--11526. \url{https://pmc.ncbi.nlm.nih.gov/articles/PMC1183554/} Bekhbat M, Li Z, Mehta ND, et al.~(2022). Functional connectivity in reward circuitry and symptoms of anhedonia as therapeutic targets in depression with high inflammation: Evidence from a dopamine challenge study. \emph{Molecular Psychiatry} 27:4113--4121. \url{https://www.nature.com/articles/s41380-022-01715-3} Berridge KC, Robinson TE (1998). What is the role of dopamine in reward: hedonic impact, reward learning, or incentive salience? \emph{Brain Research Reviews} 28:309--369. \url{https://www.sciencedirect.com/science/article/abs/pii/S0165017398000195} Block ML, Zecca L, Hong JS (2007). Microglia-mediated neurotoxicity: uncovering the molecular mechanisms. \emph{Nature Reviews Neuroscience} 8:57--69. \url{https://www.nature.com/articles/nrn2038} Dantzer R, O'Connor JC, Freund GG, Johnson RW, Kelley KW (2008). From inflammation to sickness and depression: when the immune system subjugates the brain. \emph{Nature Reviews Neuroscience} 9:46--56. \url{https://www.nature.com/articles/nrn2297} Gard DE, Gard MG, Kring AM, John OP (2006). Anticipatory and consummatory components of the experience of pleasure: A scale development study. \emph{Journal of Abnormal Psychology} 115:817--825. \url{https://pubmed.ncbi.nlm.nih.gov/17100538/} Gerfen CR, Surmeier DJ (2011). Modulation of striatal projection systems by dopamine. \emph{Annual Review of Neuroscience} 34:441--466. \url{https://www.annualreviews.org/content/journals/10.1146/annurev-neuro-061010-113641} He DY, McGough NNH, Ravindranathan A, et al.~(2005). Glial cell line-derived neurotrophic factor mediates the desirable actions of the anti-addiction drug ibogaine against alcohol consumption. \emph{Journal of Neuroscience} 25:619--628. \url{https://www.jneurosci.org/content/25/3/619} He DY, Ron D (2006). Autoregulation of glial cell line-derived neurotrophic factor expression: Implications for the long-lasting actions of the anti-addiction drug, Ibogaine. \emph{FASEB Journal} 20:E1820--E1827. \url{https://faseb.onlinelibrary.wiley.com/doi/10.1096/fj.06-6394fje} Heydorn A, Søndergaard BP, Ersbøll B, et al.~(2004). A library of 7TM receptor C- terminal tails: Interactions with the proposed post-endocytic sorting proteins ERM- binding phosphoprotein 50 (EBP50), N-ethylmaleimide-sensitive factor (NSF), sorting nexin 1 (SNX1), and G protein-coupled receptor-associated sorting protein (GASP). \emph{Journal of Biological Chemistry} 279:44989--44998. \url{https://pubmed.ncbi.nlm.nih.gov/15304494/} Hoy SM (2022). Eladocagene exuparvovec: First approval. \emph{Molecular Diagnosis \& Therapy} 26:603--610. \url{https://link.springer.com/article/10.1007/s40291-022-00608-7} Indave BI, Minozzi S, Pani PP, Amato L (2016). Antipsychotic medications for cocaine dependence. \emph{Cochrane Database of Systematic Reviews}, CD006306. \url{https://www.cochranelibrary.com/cdsr/doi/10.1002/14651858.CD006306.pub3/full} Jha MK, Minhajuddin A, Gadad BS, Trivedi MH (2017). Can C-reactive protein inform antidepressant medication selection in depressed outpatients? Findings from the CO- MED trial. \emph{Psychological Medicine} 47:1136--1148. \url{https://pubmed.ncbi.nlm.nih.gov/28090845/} McClarty B, Rodriguez G, Dong H (2021). Dose effects of histone deacetylase inhibitor tacedinaline (CI-994) on antipsychotic haloperidol-induced motor and memory side effects in aged mice. \emph{Frontiers in Neuroscience} 15:674745. \url{https://www.frontiersin.org/journals/neuroscience/articles/10.3389/fnins.2021.674745/} full Krasnova IN, Cadet JL (2009). Methamphetamine toxicity and messengers of death. \emph{Neuropharmacology} 56(Suppl 1):83--89. \url{https://pmc.ncbi.nlm.nih.gov/articles/PMC2631950/} Lohr KM, Masoud ST, Salahpour A, Miller GW (2017). Membrane transporters as mediators of synaptic dopamine dynamics: Implications for disease. \emph{Neurochemical Research} 42:2639--2662. \url{https://link.springer.com/article/10.1007/s11064-017-2259-4} Mizoguchi K, Yuzurihara M, Ishige A, Sasaki H, Chui DH, Tabira T (2000). Chronic stress induces impairment of spatial working memory because of prefrontal dopaminergic dysfunction. \emph{Neuroscience} 101:281--288. \url{https://pubmed.ncbi.nlm.nih.gov/11074152/} Moore LD, Le T, Fan G (2013). DNA methylation and its basic function. \emph{Neuropsychopharmacology} 38:23--38. \url{https://www.nature.com/articles/npp2012112} Montalvo-Ortiz JL, Fisher DW, Rodríguez G, Fang D, Csernansky JG, Dong H (2017). Histone deacetylase inhibitors reverse age-related increases in side effects of haloperidol in mice. \emph{Psychopharmacology} 234:2385--2398. \url{https://pubmed.ncbi.nlm.nih.gov/28421257/} Nestler EJ (2008). Transcriptional mechanisms of addiction: Role of ΔFosB. \emph{Nature Reviews Neuroscience} 9:844--858. \url{https://www.nature.com/articles/nrn2491} Nohesara S, Ghadirivasfi M, Barati M, et al.~(2016). Methamphetamine-induced psychosis is associated with DNA hypomethylation and increased expression of AKT1 and key dopaminergic genes. \emph{American Journal of Medical Genetics Part B: Neuropsychiatric Genetics} 171:1180--1189. \url{https://pmc.ncbi.nlm.nih.gov/articles/PMC7115129/} Renthal W, Maze I, Krishnan V, et al.~(2007). Histone deacetylase 5 epigenetically controls behavioral adaptations to chronic emotional stimuli. \emph{Neuron} 56:517--529. \url{https://pubmed.ncbi.nlm.nih.gov/17988634/} Renthal W, Carle TL, Maze I, et al.~(2008). ΔFosB mediates epigenetic desensitization of the c-fos gene after chronic amphetamine exposure. \emph{Journal of Neuroscience} 28:7344--7349. \url{https://www.jneurosci.org/content/28/29/7344} Renthal W, Kumar A, Xiao G, et al.~(2009). Genome-wide analysis of chromatin regulation by cocaine reveals a role for sirtuins. \emph{Neuron} 62:335--348. \url{https://pubmed.ncbi.nlm.nih.gov/19447090/} Robertson CL, Ishibashi K, Chudzynski J, et al.~(2016). Effect of exercise training on striatal dopamine D2/D3 receptors in methamphetamine users during behavioral treatment. \emph{Neuropsychopharmacology} 41:1629--1636. \url{https://pmc.ncbi.nlm.nih.gov/articles/PMC4832026/} Schenberg EE, de Castro Comis MA, Chaves BR, da Silveira DX (2014). Treating drug dependence with the aid of ibogaine: A retrospective study. \emph{Journal of Psychopharmacology} 28:993--1000. \url{https://pubmed.ncbi.nlm.nih.gov/25271214/} Schultz W (2015). Neuronal reward and decision signals: From theories to data. \emph{Neuron} 86:646--664. \url{https://www.cell.com/neuron/fulltext/S0896-6273(15)00323-1} Snaith RP, Hamilton M, Morley S, Humayan A, Hargreaves D, Trigwell P (1995). A scale for the assessment of hedonic tone: The Snaith-Hamilton Pleasure Scale. \emph{British Journal of Psychiatry} 167:99--103. \url{https://pubmed.ncbi.nlm.nih.gov/7551619/} Thanos PK, Volkow ND, Freimuth P, et al.~(2001). Overexpression of dopamine D2 receptors reduces alcohol self-administration. \emph{Journal of Neurochemistry} 78:1094-- 1103. \url{https://pubmed.ncbi.nlm.nih.gov/11553683/} Thanos PK, Rivera SN, Weaver K, et al.~(2005). Dopamine D2R DNA transfer into the nucleus accumbens attenuates cocaine self-administration in rats. \emph{Synapse} 54:119-- 128. \url{https://pubmed.ncbi.nlm.nih.gov/15352133/} Treadway MT, Buckholtz JW, Schwartzman AN, Lambert WE, Zald DH (2009). Worth the `EEfRT'? The effort expenditure for rewards task as an objective measure of motivation and anhedonia. \emph{PLoS ONE} 4:e6598. \url{https://journals.plos.org/plosone/article?id=10.1371/journal.pone.0006598} Treadway MT, Zald DH (2011). Reconsidering anhedonia in depression: Lessons from translational neuroscience. \emph{Neuroscience \& Biobehavioral Reviews} 35:537--555. \url{https://pubmed.ncbi.nlm.nih.gov/20603146/} Treadway MT, Buckholtz JW, Cowan RL, et al.~(2012). Dopaminergic mechanisms of individual differences in human effort-based decision-making. \emph{Journal of Neuroscience} 32:6170--6176. \url{https://www.jneurosci.org/content/32/18/6170} Uher R, Tansey KE, Dew T, et al.~(2014). An inflammatory biomarker as a differential predictor of outcome of depression treatment with escitalopram and nortriptyline. \emph{American Journal of Psychiatry} 171:1278--1286. \url{https://psychiatryonline.org/doi/10.1176/appi.ajp.2014.14010094} Vickery RG, von Zastrow M (1999). Distinct dynamin-dependent and -independent mechanisms target structurally homologous dopamine receptors to different endocytic membranes. \emph{Journal of Cell Biology} 144:31--43. \url{https://pubmed.ncbi.nlm.nih.gov/9885242/} Volkow ND, Chang L, Wang GJ, et al.~(2001). Loss of dopamine transporters in methamphetamine abusers recovers with protracted abstinence. \emph{Journal of Neuroscience} 21:9414--9418. \url{https://pubmed.ncbi.nlm.nih.gov/11717374/} Volkow ND, Wang GJ, Smith L, et al.~(2015). Recovery of dopamine transporters with methamphetamine detoxification is not linked to changes in dopamine release. \emph{NeuroImage} 121:20--28. \url{https://pubmed.ncbi.nlm.nih.gov/26208874/} Volkow ND, Koob GF, McLellan AT (2016). Neurobiologic advances from the brain disease model of addiction. \emph{New England Journal of Medicine} 374:363--371. \url{https://www.nejm.org/doi/full/10.1056/NEJMra1511480} Volkow ND, Wise RA, Baler R (2017). The dopamine motive system: Implications for drug and food addiction. \emph{Nature Reviews Neuroscience} 18:741--752. \url{https://www.nature.com/articles/nrn.2017.130} Wang GJ, Volkow ND, Chang L, et al.~(2004). Partial recovery of brain metabolism in methamphetamine abusers after protracted abstinence. \emph{American Journal of Psychiatry} 161:242--248. \url{https://psychiatryonline.org/doi/10.1176/appi.ajp.161.2.242} Wang GJ, Smith L, Volkow ND, et al.~(2012). Decreased dopamine activity predicts relapse in methamphetamine abusers. \emph{Molecular Psychiatry} 17:918--925. \url{https://pmc.ncbi.nlm.nih.gov/articles/PMC3261322/} Wernicke C, Hellmann J, Zięba B, et al.~(2010). 9-Methyl-β-carboline has restorative effects in an animal model of Parkinson's disease. \emph{Pharmacological Reports} 62:35-- 53. \url{https://pubmed.ncbi.nlm.nih.gov/20360614/} Wohleb ES, Hanke ML, Corona AW, et al.~(2011). β-Adrenergic receptor antagonism prevents anxiety-like behavior and microglial reactivity induced by repeated social defeat. \emph{Journal of Neuroscience} 31:6277--6288. \url{https://www.jneurosci.org/content/31/17/6277} Yuen EY, Liu W, Karatsoreos IN, Feng J, McEwen BS, Yan Z (2009). Acute stress enhances glutamatergic transmission in prefrontal cortex and facilitates working memory. \emph{Proceedings of the National Academy of Sciences} 106:14075--14079. \url{https://www.pnas.org/doi/10.1073/pnas.0906791106} Yuen EY, Wei J, Liu W, Zhong P, Li X, Yan Z (2012). Repeated stress causes cognitive impairment by suppressing glutamate receptor expression and function in prefrontal cortex. \emph{Neuron} 73:962--977. \url{https://www.cell.com/neuron/fulltext/S0896-} 6273(12)00087-5
